\textbf{5.} Sean $q:X\to Y$ y $r:Y\to Z$ funciones recubridoras y considere $p=r\circ q$. Demuestre que si $r^{-1}(z)$ es finito para cada $z\in Z$, entonces $p$ es una función recubridora.
\begin{proof}
   Supongamos que $r^{-1}(z)$ es finito para cada $z\in Z$.\\\\
   Tenemos que $p$ es continua y sobreyectiva por ser composición de funciones continuas y sobreyectivas.\\\\
   Sea $z\in Z$ cualquiera. Como $r:Y\to Z$ es una función recubridora, existen $U\subseteqab Z$ con $z\in U$ y $\left( V_i\right)_{i\in I}$ una colección de abiertos disyuntos de $Y$ tal que $\bigsqcup_{i\in I}V_i=r^{-1}(U)$ y $\left.r\right|_{V_i}$ es un homeomorfismo de $V_i$ en $U$ para cada $i\in I$. Notemos que para cada $i\in I$ existe $z_i\in V_i$ tal que $r(z_i)=\left.r\right|_{V_i}(z_i)=z$, de modo que $z_i\in r^{-1}(z)$. Para cualesquiera $i\ne j$ en $I$ se tiene $V_i\cap V_j=\emptyset$, lo que implica $z_i\ne z_j$, con lo cual para cada $i\in I$ hay un $z_i$ distinto en $r^{-1}(z)$. Como $r^{-1}(z)$ es finito, se sigue que $I$ es finito.\\\\
   Sea $i\in I$ cualquiera. Como $z_i\in V\subseteq Y$, entonces $z_i\in Y$ y como $q:X\to Y$ es una función recubridora, existen $W_i\subseteqab Y$ con $z_i\in W_i$ y $\left( A_{i,k}\right)_{k\in I_i}$ una colección de abiertos disyuntos de $X$ tales que $\bigsqcup_{k\in I_i}A_{i,k}=q^{-1}\left( W_i\right)$ y $\left.q\right|_{A_{i,k}}$ es un homeomorfismo de $A_{i,k}$ en $W_i$ para cada $k\in I_i$. Tomemos $B_i:=W_i\cap V_i\subseteqab Y$. Tenemos en particular $B_i\subseteqab V_i$, luego $\left.r\right|_{V_i}\left( B_i\right)\subseteqab U\subseteqab Z$, luego $\left.r\right|_{V_i}\left( B_i\right)\subseteqab Z$.\\\\
   Tomemos $\widetilde{U}:=\bigcap_{i\in I}\left.r\right|_{V_i}\left( B_i\right)$. Como $I$ es finito, $\widetilde{U}\subseteqab Z$. Dado $i\in I$ cualquiera se tiene $z_i\in W_i\cap V_i=B_i$, luego $z=r(z_i)=\left.r\right|_{V_i}(z_i)\in \left.r\right|_{V_i}\left( B_i\right)$, de modo que $z_i\in\widetilde{U}$; además $\left.r\right|_{V_i}\left( B_i\right)\subseteq U$, luego $\widetilde{U}\subseteq U$. Como $\widetilde{U}\subseteqab Z$ entonces $\widetilde{U}\subseteqab U$.\\\\
Sea $i\in I$ cualquiera. Tomemos $\widetilde{B}_i:=\left.r\right|_{V_i}^{-1}\left( \widetilde{U}\right)\subseteqab V_i$. Como $V_i\subseteqab Y$ entonces $\widetilde{B}_i\subseteqab Y$. En particular $\widetilde{U}\subseteq \left.r\right|_{V_i}\left( B_i\right)$, luego $\widetilde{B}_i=\left.r\right|_{V_i}^{-1}\left( \widetilde{U}\right)\subseteq \left.r\right|_{V_i}^{-1}\left( \left.r\right|_{V_i} \left( B_i\right)\right)=B_i\subseteq W_i$, luego $\widetilde{B}_i\subseteqab W_i$. Sea $k\in I_i$ cualquiera. Tomemos $\widetilde{A}_{i,k}:=\left.q\right|_{A_{i,k}}^{-1}\left( \widetilde{B}_i\right)\subseteqab A_{i,k}$. Como $A_{i,k}\subseteqab X$ entonces $\widetilde{A}_{i,k}\subseteqab X$.\\\\
   Para cada $i\in I$ sabemos que $\left( A_{i,k}\right)_{k\in I_i}$ es una colección de conjuntos disyuntos y $\widetilde{A}_{i,k}\subseteq A_{i,k}$ para todo $k\in I_i$, luego $\left( \widetilde{A}_{i,k}\right)_{k\in I_i}$ es una colección de conjuntos disyuntos. Sean $i_1,i_2\in I$ cualesquiera con $i_1\ne i_2$, y tomemos $k_1\in I_{i_1}$ y $k_2\in I_{i_2}$ cualesquiera. Razonando por contradicción supongamos que existe $x\in\widetilde{A}_{i_1,k_1}\cap\widetilde{A}_{i_2,k_2}$. Como $x\in\widetilde{A}_{i_1,k_1}$ entonces $\left.q\right|_{A_{i_1,k_1}}(x)\in\widetilde{B}_{i_1}\subseteq B_{i_1}\subseteq V_{i_1}$. Análogamente $x\in V_{i_2}$, de modo que $x\in V_{i_1}\cap V_{i_2}$, contradiciendo que $V_{i_1}$ y $V_{i_2}$ son disyuntos. Por tanto $\widetilde{A}_{i_1,k_1}\cap\widetilde{A}_{i_2,k_2}=\emptyset$. Esto implica que la colección $\left( \widetilde{A}_{i,k}\right)_{i\in I, k\in I_i}$ de subconjuntos abiertos de $X$ es disyunta.\\\\
   Veamos que $p^{-1}\left( \widetilde{U}\right)=\bigsqcup_{i\in I, k\in I_i}\widetilde{A}_{i,k}$.\\\\
   Sea $x\in p^{-1}\left( \widetilde{U}\right)$ cualquiera. Se tiene que $p(x)\in\widetilde{U}$. En particular $U\ni p(x)=r(q(x))$, de modo que $q(x)\in p^{-1}(U)=\bigsqcup_{i\in I}V_i$, luego $q(x)\in V_i$ para algún $i\in I$. Así, $\left.r\right|_{V_i}(q(x))=r(q(x))=p(x)\in \widetilde{U}$, luego $q(x)\in \left.r\right|_{V_i}^{-1}\left( \widetilde{U}\right)=\widetilde{B}_i\subseteq B_i\subseteq W_i$, con lo cual $x\in q^{-1}\left( W_i\right)=\bigsqcup_{k\in I_i}A_{i,k}$, de modo que existe $k\in I_i$ tal que $x\in A_{i,k}$. Notemos que $\left.q\right|_{A_{i,k}}(x)=q(x)\in\widetilde{B}_i$, luego $x\in \left.q\right|_{A_{i,k}}^{-1}\left( \widetilde{B}_i\right)=\widetilde{A}_{i,k}$. Con esto $x\in \bigsqcup_{i\in I, k\in I_i}\widetilde{A}_{i,k}$ y por tanto $p^{-1}\left( \widetilde{U}\right)\subseteq \bigsqcup_{i\in I, k\in I_i}\widetilde{A}_{i,k}$.\\\\\
   Sea $x\in \bigsqcup_{i\in I, k\in I_i}\widetilde{A}_{i,k}$. Existen $i\in I$ y $k\in I_i$ tales que $x\in \widetilde{A}_{i,k}=\left.q\right|_{A_{i,k}}^{-1}\left( \widetilde{B}_i\right)$, luego $q(x)=\left.q\right|_{A_{i,k}}(x)\in\widetilde{B}_i=\left.r\right|_{V_i}^{-1}\left( \widetilde{U}\right)$, es decir $p(x)=r(q(x))=\left.r\right|_{V_i}(q(x))\in \widetilde{U}$, con lo cual $x\in p^{-1}\left( \widetilde{U}\right)$. Así, $\bigsqcup_{i\in I, k\in I_i}\widetilde{A}_{i,k}\subseteq p^{-1}\left( \widetilde{U}\right)$.\\\\
   De lo anterior, $p^{-1}\left( \widetilde{U}\right)=\bigsqcup_{i\in I, k\in I_i}\widetilde{A}_{i,k}$.\\\\
   Sean $i\in I$ y $k\in I_i$ cualesquiera. Veamos que $\left.p\right|_{\widetilde{A}_{i,k}}$ es un homeomorfismo de $\widetilde{A}_{i,k}$ en $\widetilde{U}$.\\\\
   Notemos que $\left.q\right|_{A_{i,k}}$ es un homeomorfismo de $A_{i,k}$ en $W_i$ y $\widetilde{A}_{i,k}=\left.q\right|_{A_{i,k}}^{-1}\left( \widetilde{B}_i\right)$; como $\left.q\right|_{A_{i,k}}$ es en particular biyectiva, tenemos $\left.q\right|_{A_{i,k}}\left( \widetilde{A}_{i,k}\right)=\left( \widetilde{B}_i\right)$, por tanto $\left.\left( \left.q\right|_{A_{i,k}}\right)\right|_{\widetilde{A}_{i,k}}=\left.q\right|_{\widetilde{A}_{i,k}}$ es un homeomorfismo de $\widetilde{A}_{i,k}$ en $\widetilde{B}_i$. Igualmente $\left.r\right|_{V_i}$ es un homeomorfismo de $V_i$ en $U$ y $\widetilde{B}_i=\left.r\right|_{V_i}^{-1}\left( \widetilde{U}\right)$; como $\left.r\right|_{V_i}$ es en particular biyectiva, tenemos $\left.r\right|_{V_i}\left( \widetilde{B}_i\right)=\widetilde{U}$, luego $\left.\left( \left.r\right|_{V_i}\right)\right|_{\widetilde{B}_{i}}=\left.r\right|_{\widetilde{B}_{i}}$ es un homeomorfismo de $\widetilde{B}_i$ en $\widetilde{U}$. Por tanto $\left.r\right|_{\widetilde{B}_i}\circ\left.q\right|_{\widetilde{A}_{i,k}}$ es un homeomorfismo de $\widetilde{A}_{i,k}$ en $\widetilde{U}$.\\\\
   Notemos que $\left.p\right|_{\widetilde{A}_{i,k}}$ y $\left.r\right|_{\widetilde{B}_i}\circ\left.q\right|_{\widetilde{A}_{i,k}}$ tienen el mismo dominio $\widetilde{A}_{i,k}$, y para cada $x\in\widetilde{A}_{i,k}$ se tiene
   $$
   \begin{aligned}
      \left.p\right|_{\widetilde{A}_{i,k}}(x) &= p(x)\\
                                              &= r(q(x))\\
                                              &= \left.r\right|_{\widetilde{B}_i}\left( \left.q\right|_{\widetilde{A}_{i,k}}(x)\right)\\
                                              &= \left( \left.r\right|_{\widetilde{B}_i}\circ\left.q\right|_{\widetilde{A}_{i,k}}\right)(x).
   \end{aligned}
   $$
   Como $\left.r\right|_{\widetilde{B}_i}\circ\left.q\right|_{\widetilde{A}_{i,k}}$ es un homeomorfismo de $\widetilde{A}_{i,k}$ en $\widetilde{U}$, obtenemos que $\left.p\right|_{\widetilde{A}_{i,k}}$ es un homeomorfismo de $\widetilde{A}_{i,k}$ en $\widetilde{U}$.\\\\
   Concluimos que $p$ es una función recubridora.
\end{proof}
